% =====================================================================
\chapter*{Instrucciones}
\addcontentsline{toc}{chapter}{Instrucciones}
\markboth{Instrucciones}{}
% =====================================================================

Este examen evalúa la fase~0 del curso, los preliminares matemáticos: lógica y técnicas de demostración (módulo~0.1), conjuntos, funciones y relaciones (módulo~0.2), números, álgebra y sumatorios (módulo~0.3), trigonometría y números complejos (módulo~0.4), cálculo diferencial de una variable (módulo~0.5), cálculo integral y series (módulo~0.6), vectores y funciones de varias variables (módulo~0.7), combinatoria y probabilidad (módulo~0.8) y estadística (módulo~0.9). Se valoran por igual el cálculo correcto y la demostración bien redactada: en los apartados que piden demostrar algo, cada paso debe estar justificado por una definición, un teorema de los apuntes o una cuenta explícita.

\section*{Condiciones}
\begin{center}
\small
\begin{tabularx}{\textwidth}{@{}lX@{}}
\toprule
Duración & 4 horas. \\
Material permitido & Calculadora científica no programable. No se permiten apuntes, libros, formularios propios ni dispositivos conectados. Los valores de las tablas estadísticas necesarios figuran en el \cref{tab:datos}. \\
Puntuación & 100 puntos: Parte~A, 40 puntos (20 cuestiones de 2 puntos); Parte~B, 60 puntos (9 problemas: seis de 7 puntos y tres de 6 puntos). \\
Aprobado & 70 puntos (70\,\%) o más. \\
Módulos & 0.1 a 0.9 (véase el reparto por módulo en el \cref{ap:autoevaluacion}). \\
Proyecto de fase & Se evalúa aparte sobre 100 puntos (\cref{cap:proyecto}) y se entrega en el plazo que fije el profesor. \\
\bottomrule
\end{tabularx}
\end{center}

\section*{Normas}
\begin{enumerate}[label=\arabic*.]
  \item En la Parte~A hay una sola opción correcta por cuestión. Las respuestas erróneas no restan.
  \item En la Parte~B debe mostrarse el razonamiento completo. Una respuesta numérica sin desarrollo no puntúa. Los resultados exactos se dejan en forma exacta (fracciones, radicales, $\pi$, $e$, $\ln 2$) y, cuando se pida un valor decimal, se dan cuatro cifras significativas.
  \item En las demostraciones puede usarse sin probarlo cualquier resultado enunciado en los apuntes de la fase~0, citándolo por su nombre (por ejemplo, «por el teorema del valor medio»), salvo que el enunciado pida precisamente demostrar ese resultado.
  \item Los ángulos se miden en radianes salvo que se indique otra cosa. $\ln$ es el logaritmo natural y $\log_2$, el logaritmo en base~2.
  \item La sucesión de Fibonacci se define por $F_0=0$, $F_1=1$ y $F_{n+2}=F_{n+1}+F_n$ para $n\ge0$; con esta definición, $F_2=1$, como en los apuntes del módulo~0.1.
\end{enumerate}

\section*{Datos de referencia}
\begin{table}[h]
\centering\small
\caption{Valores de referencia del examen.}\label{tab:datos}
\begin{tabularx}{\textwidth}{@{}lX@{}}
\toprule
Dato & Valor \\
\midrule
Constantes & $\ln 2=0{,}6931$; $\ln 5=1{,}6094$; $\ln 10=2{,}3026$; $e=2{,}7183$; $e^{2}=7{,}3891$; $e^{-4}=0{,}0183$; $\pi=3{,}1416$ \\
Normal estándar, función de distribución & $\Phi(1{,}645)=0{,}9500$; $\Phi(1{,}90)=0{,}9713$; $\Phi(1{,}96)=0{,}9750$ \\
Cuantiles normales & $z_{0{,}95}=1{,}645$; $z_{0{,}975}=1{,}960$ \\
Cuantiles de la $t$ de Student & $t_{7;0{,}975}=2{,}365$; $t_{8;0{,}975}=2{,}306$ \\
Raíces & $\sqrt2=1{,}4142$; $\sqrt3=1{,}7321$; $\sqrt6=2{,}4495$; $\sqrt8=2{,}8284$; $\sqrt{148}=12{,}166$ \\
\bottomrule
\end{tabularx}
\end{table}
